\section{Fast Apply: convertir la salida del modelo en ediciones verificables}

Después de que el modelo genera un patch, el sistema todavía tiene que localizar la posición en el archivo, manejar las modificaciones simultáneas del usuario, aplicar el diff, conservar el formato y disparar los diagnostics. Los materiales públicos que Cursor publicó más tarde sobre Fast Apply muestran que un specialized model puede encargarse de la edit application de alto rendimiento; esta sección lo trata como un durable design pattern, sin trasladar retroactivamente a la clase el desempeño posterior.

\subsection{Plan, generate, apply, validate}

La sección anterior resolvió «de dónde obtener la salida del modelo»; esta sección resuelve «cómo hacer que entre de forma segura al editor state». Al leer la figura, conviene fijarse en la separación de responsabilidades entre plan y apply, y en por qué validation/undo son infraestructura del producto. Si apply modifica en silencio la posición equivocada, se destruye la confianza aunque el texto sea de alta calidad.

\lecturefigure{13-fast-apply.png}{El planning model y el specialized apply model se reparten la comprensión y la edición a alta velocidad.}{Subtítulos locales 04:20--05:20; materiales oficiales de Cursor Fast Apply.}

\begin{knowledgebox}{Lectura de la figura: los cuatro frentes de falla de la patch application}
Errores de localización, context drift, fallas de sintaxis/tipos y ediciones concurrentes del usuario. El sistema debe usar precondition de file/version, structured edit, parser/diagnostic, visible diff y atomic undo. Para cambios grandes, también se necesita aplicar por partes (apply por lotes) con validation intermedia, para evitar que un único patch gigantesco bloquee el editor.
\end{knowledgebox}

\begin{importantbox}{La velocidad de interacción no es solo tokens/s}
El tiempo que percibe el usuario incluye el first useful preview, la edit stabilization y el validation complete. Un apply path más lento, pero que puede mostrarse de forma progresiva, tiene pocos conflictos y permite un undo rápido, puede ser mejor que uno con mayor token throughput puro. Las métricas deben incluir accepted edit, rework y rollback, no solo la velocidad de generación.
\end{importantbox}

\subsection{Resumen de la sección}

Fast Apply convierte AI coding de un problema de generación de texto en un problema de state transition. El modelo se encarga de proponer las modificaciones; el sistema debe encargarse de precondition, application, validation, visibility y recovery.
